% RESEARCH_DRAFT_SHA256: a43f47b2a856109f5860bea1cd7fae4c6774fd00ad9ca33e8a2577235e017f35
% Layout and source-reference wrappers expanded for the web renderer.
\documentclass[11pt]{article}
\usepackage[margin=25mm]{geometry}
\usepackage{fontspec,unicode-math}
\setmainfont{Latin Modern Roman}
\setmathfont{Latin Modern Math}
\usepackage{amsmath,amssymb,mathtools,xurl,hyperref}
\hypersetup{colorlinks=true,urlcolor=blue,linkcolor=blue}
\setcounter{secnumdepth}{0}
\setlength{\parindent}{0pt}
\setlength{\parskip}{5pt}
\setlength{\emergencystretch}{3em}
\newcommand{\sref}[2]{\hyperlink{source:#1:#2}{[S#2]}}
\newcommand{\cataloguescope}[1]{\paragraph{Source record.}#1}
\newcommand{\floor}[1]{\left\lfloor#1\right\rfloor}
\begin{document}
\section*{Research-attempt provenance}
This standalone research attempt retains the complete original definition body below,
including its historical source and preparation notes. Those notes describe the
earlier catalogue preparation; the new attempt and its verification are identified
separately after that body. No original ID, catalogue status, or source statement is changed.

% UnsolvedMath ID 30005086; source catalogue code OWR-10252928-001
\section{Unspecified Conic Optimization Research Problem}\label{30005086}
{\small UnsolvedMath ID 30005086 · Computer Science}
\subsection{Definitions and mathematical statement}

\paragraph{Shared notation.}Standard number systems are $\mathbb N,\mathbb Z,\mathbb Q,\mathbb R,\mathbb C$. Specialized notation and index ranges follow the cited source and the definitions below.


Conic optimization minimizes a linear functional over the intersection of a convex cone and an affine subspace. A cone is closed under nonnegative scaling; different cones yield different optimization classes.
\paragraph{Record scope and qualifications.}
The catalogue title is retained as metadata. Named certificates, contracts, bounds or frameworks in that title are not asserted as results of the original source extract.
The supplied registry flags this extraction as unrecoverable. It is displayed as a literal source fragment; the definitions provide context rather than a replacement formal statement.

This is an abstract or workshop introduction, not a complete individual problem statement.

Source-record limitation: Source flags the statement as unrecoverable: The official report shows that this is workshop-level introductory prose rather than a standalone conic-optimization problem. No separate clean mathematical statement is supplied; this source ID preserves its provenance.
The source text below is rendered literally where necessary; damaged formulas and composite extracts are not silently repaired into a new mathematical claim.
The supplied catalogue marks this record as unrecoverable. Its extract and terminology are documented, but a complete original problem has not been reconstructed.
\paragraph{Statement recorded in the supplied source.}
Abstract. From a mathematical perspective, optimization is the science of proving inequalities. In this sense, computational optimization is a method for computer-assisted proofs. Conic (linear) optimization is the problem of minimizing a linear func- tional over the intersection of a convex cone with an affine subspace of a topological vector space. For many cones this problem is computationally tractable, and as a result there is a growing number of computer-assisted proofs using conic optimization in discrete geometry, (extremal) graph the- ory, numerical analysis, and other fields, the most famous example perhaps being the proof of the Kepler Conjecture. The aim of this workshop was to bring researchers from these diverse fields together to work towards expanding the current scope of conic optimization as a method of generating proofs, and to identify problems and challenges to work on together. Mathematics Subject Classification (2020): 90Cxx, 90C22, 52-xx, 14Pxx, 65-xx, 05Cxx. \href{https://doi.org/10.4171/owr/2022/20}{[S2]}
\paragraph{Source status.}
The supplied catalogue labels this source record open; that metadata was checked against the supplied snapshot on 2026-10-11. The label is preserved for provenance and is not an independent verification of present open status or mathematical correctness. The catalogue literature assessment is dated 2026-08-22.
\subsection{Short English statement}
The fragment introduces a workshop on conic optimization and computer-assisted proofs.
\subsection{Sources}
\begin{itemize}
\item[S1] Ulam AI and the UnsolvedMath Contributors, UnsolvedMath: A Curated Collection of Open Mathematics Problems, record 30005086 (OWR-10252928-001). Catalogue curation and metadata: CC BY 4.0. The source extract is retained; definitions, English summaries and qualifications are editorial additions.
\url{https://huggingface.co/datasets/ulamai/UnsolvedMath}.
\item[S2] Conic Linear Optimization for Computer-Assisted Proofs (2023), Oberwolfach Reports, DOI 10.4171/owr/2022/20, pp. 1039--1089
\url{https://doi.org/10.4171/owr/2022/20}.
\item Source-specific Oberwolfach Reports licences, including CC BY-SA 4.0 where applicable, are retained; see THIRD\_PARTY\_NOTICES.md.
\end{itemize}
\paragraph{Preparation and review.}
The original source and UnsolvedMath contributors retain attribution for the source extract. Codex (GPT-6; exact serving revision not exposed) assisted with definitions, English summaries, scope notes and formatting. No new human mathematical authorship or independently verified solution is asserted. Permanent source ID 30005086 and code OWR-10252928-001 checked against the complete supplied registry and current repository. Distinct upstream records, including repeated questions and broad or damaged extracts, are retained for complete coverage of the source CS category. A repeated record is not a new independent conjecture. Complete mathematical-equivalence and current-literature review has not been performed. Preparation and final review follow contribution prompts C and D, with source gaps disclosed. Structural and rendering checks do not establish a new theorem or open status.
\label{end:30005086}

\section{Literature and context check}

The primary report's opening page 1039 is a workshop abstract, exactly as the
preserved extract suggests. It is not a single conjecture with quantified input
or output. I read that page and the organizers' discussion of exact SDP bounds
on page 1040. The following is an explicitly auxiliary finite-dimensional
certificate calculation, relevant to the computer-assisted-proof theme. It does
not turn the workshop agenda into a theorem purportedly solved here.

\section{Research attempt: Residual-aware semidefinite proof certification}

\paragraph{A concrete certification route.}
Consider real symmetric matrices and the optimization problem
\[
 p_* =\inf\{\langle C,X\rangle:\mathcal A(X)=b,\ X\succeq0\},
 \qquad \langle C,X\rangle=\operatorname{tr}(CX).
\]
Suppose every feasible $X$ satisfies $\operatorname{tr}X\leq R$ for a known
$R\geq0$. Given a vector $y$, a genuinely positive semidefinite matrix $S$,
and a certified residual bound
\[
 E=C-\mathcal A^*y-S,\qquad \|E\|_{\rm op}\leq\varepsilon,
\]
we obtain the rigorous lower bound
\[
 p_*\geq b^Ty-\varepsilon R.
\]
Indeed, $\langle C,X\rangle=b^Ty+\langle S,X\rangle+\langle E,X\rangle$.
Diagonalize $X=\sum_i\lambda_i u_i u_i^T$ with $\lambda_i\geq0$.
Then $\langle S,X\rangle=\sum_i\lambda_i u_i^TSu_i\geq0$ and
$\langle E,X\rangle\geq-\varepsilon\sum_i\lambda_i\geq-\varepsilon R$.
No dual attainment or Slater condition is needed for this elementary weak-duality
certificate. With an exactly feasible candidate $X_0$, the same calculation gives
a certified interval
$[b^Ty-\varepsilon R,\langle C,X_0\rangle]$ containing $p_*$.
A rational factorization $S=LL^T$ is one possible exact PSD certificate; verifying
its entries suffices because $z^TSz=\|L^Tz\|^2\geq0$.

\paragraph{Why small floating residual alone fails.}
In dimension one, let $C=0$, impose $X=1/\varepsilon$ with rational
$\varepsilon>0$, and take $y=\varepsilon$, $S=0$.
The true optimum is zero, while the uncorrected proposed lower bound is
$b y=1$. Its dual residual is only $E=-\varepsilon$.
The trace factor is $R=1/\varepsilon$, so the corrected bound is exactly zero.
As $\varepsilon\downarrow0$, the residual tends to zero but the false lower-bound
error stays one. Therefore no error guarantee depending only on residual size
and tending to zero can hold uniformly without a feasible-scale bound.
If the feasible set is unbounded, even the one-dimensional problem
$\inf_{X\geq0}(-\varepsilon X)=-\infty$ has the approximate certificate
$y=0,S=0,E=-\varepsilon$. A small residual cannot certify a finite lower bound.

\paragraph{PSD rounding must also be certified.}
The rank-one PSD matrix $\left(\begin{smallmatrix}1&1\\1&1\end{smallmatrix}\right)$
becomes
\[
 S_\delta=\begin{pmatrix}1&1\\1&1-\delta\end{pmatrix}
\]
after an entry perturbation of size $\delta>0$.
The vector $(1,-1)$ has quadratic form $-\delta$, so $S_\delta$ is not PSD,
however small the perturbation. One valid repair is to establish a certified
$S_\delta\succeq-\eta I$ and then use $S_\delta+\eta I\succeq0$, charging
the resulting residual $\eta I$ to the same trace bound. Such a repair must
be proved for the particular numerical output; proximity to a floating PSD
matrix is not an exact certificate.

\section{Verification and adversarial audit}

The residual sign was checked on the scalar false-certificate example. All bounds
are inequalities for every feasible point, so neither existence of an optimizer
nor boundedness of the dual feasible set is being assumed. The interval claim
requires a feasible $X_0$; without one it is only a lower bound. The rounding
counterexample supplies an explicit negative vector instead of relying on an
approximate eigenvalue calculation.

\paragraph{Audit corrections.} Explicitly required a certified trace bound and genuine PSD slack, and did
not treat a small numerical residual as a proof. The broad original abstract is
preserved as such.
\paragraph{External computation.} None; symbolic arguments and boundary checks only.
\section{Scope and remaining gaps}

The supplied record is an unspecified workshop program. The bounded-trace SDP
lemma and two obstructions are auxiliary known-style certification facts, not a
resolution of that program or a new conic optimization theorem. No concrete
source conjecture or numerical SDP instance is supplied whose optimal value
could be certified as the original target.

\paragraph{Integration review.} Codex (GPT-6; exact serving revision not exposed) reviewed the written arguments and expanded layout and source-reference wrappers. Repository metadata and the production web renderer were checked using Python 3.13 and JavaScript with MathJax 3.2.2. This model-assisted review is not independent human proof verification; no PDF compilation or formal proof checking is asserted.

\section{Completion Estimate}
\noindent\textbf{COMPLETION: 0\%}
The calculations below the preserved source statement investigate a concrete route,
but establish no substantiated new progress toward the entire original target.
Known special cases and failed sufficient conditions are not credited as a solution.
\section{References}
The original source references and catalogue attribution are preserved in the
Sources section above. The following URLs were actually consulted for this attempt;
retention of an additional original reference does not claim fresh consultation of it.
\begin{itemize}
\item \url{https://ems.press/content/serial-article-files/46957} (consulted 2026-10-11). 
\end{itemize}
\section{Model and contributor disclosure}
Codex GPT-6 completed written-proof review of this attempt at its declared scope; this is model-assisted review, not independent human proof checking.
OpenAI gpt-6-astra, requested reasoning configuration ultra, generated this research
attempt on 2026-10-11. The exact serving revision is not exposed. Codex GPT-6
performed orchestration and repository-format review; its exact serving revision
is also not exposed. The mathematical route, boundary checks and final adversarial
self-audit followed contribution prompts A then B in
\texttt{docs/contribution\_prompts.md}, after reading \texttt{CONTRIBUTING.md}.
Python 3.14.2 assembled the standalone documents; any mathematical computation is
specified in Verification. No theorem prover was used. No human mathematical
authorship, substantial human proof correction, or independent human proof checking
is asserted. Model-assisted review, source retrieval, and structural checks do not
establish correctness or novelty. Full private reasoning transcripts are unavailable
as reproducibility artifacts; the complete arguments, inputs and outcomes actually
used are stated here or in the companion batch audit.
\section{Final status}
Unresolved as a new research contribution. This is a concrete route investigation,
not a claim of new complete resolution or a current-status certification. A known
answer reproduced for context does not become a new contribution. The catalogue's
source-status label is preserved; this attempt status does not assert that a
source-labelled solved question is open.
\end{document}
